\section*{Chapitre \ref{ch:b2:asexual-reproduction} --- Reproduction asexuée et clonage chez les plantes}

\begin{solution}{exo:b2:asexual-reproduction:1}
\omterm{def:b2:asexual-reproduction:clone}{Clone}~: les descendants génétiquement identiques d'un même parent,
obtenus par mitoses. \omterm{def:b2:asexual-reproduction:clone}{Génet}~: l'individu génétique, tout ce qui est issu
d'un même zygote. \omterm{def:b2:asexual-reproduction:clone}{Ramet}~: une unité physiologiquement indépendante d'un
\omterm{def:b2:asexual-reproduction:clone}{génet}. Organes~: stolons (fraisier), rhizomes (chiendent), tubercules
(pomme de terre), bulbes (oignon), drageons (tremble), plantules
foliaires (\emph{Kalanchoe}).
\end{solution}

\begin{solution}{exo:b2:asexual-reproduction:2}
Les cambiums du greffon et du porte-greffe doivent être alignés et en
contact~; un cal comble la coupure et y différencie du xylème et du
phloème. Le porte-greffe fournit les racines (eau, minéraux, vigueur,
tolérance au sol, contrôle de la taille), le greffon la pousse
fructifère (la variété). Chaque partie garde son propre génome et les
cellules ne fusionnent jamais~: deux \omterm{def:b2:meiosis-heredity:vocabulary}{génotypes} côte à côte, soit une
chimère~; un hybride porterait un mélange des deux génomes dans chaque
cellule.
\end{solution}

\begin{solution}{exo:b2:asexual-reproduction:3}
\omterm{def:b2:asexual-reproduction:clone}{Multiplication végétative}~: une partie du corps s'enracine et donne une
nouvelle plante. \omterm{def:b2:asexual-reproduction:apomixis}{Apomixie}~: une \omterm{def:b2:angiosperm-reproduction:seed}{graine} se forme sans fécondation,
l'embryon étant un \omterm{def:b2:asexual-reproduction:clone}{clone} de la mère. \omterm{def:b2:asexual-reproduction:apomixis}{Parthénogenèse}~: un animal se
développe à partir d'un ovocyte non fécondé. Seule l'\omterm{def:b2:asexual-reproduction:apomixis}{apomixie} conserve la
\omterm{def:b2:angiosperm-reproduction:seed}{graine} --- avec son tégument, sa \omterm{def:b2:angiosperm-reproduction:seed}{dormance} et sa dispersion.
\end{solution}

\begin{solution}{exo:b2:asexual-reproduction:4}
Stériliser l'explant~; le placer sur une gélose contenant sucre,
minéraux, vitamines, et auxine et cytokinine en équilibre (cal)~;
élever le rapport cytokinine/auxine pour induire des tiges~; repiquer les
tiges toutes les quelques semaines~; les transférer sur un milieu riche
en auxine pour les enraciner~; sevrer les plantules sous forte humidité,
puis en terre.
\end{solution}

\begin{solution}{exo:b2:asexual-reproduction:5}
Chaque plante est remplacée par $1 + 24 = 25$~: au bout de 4 ans,
$25^{4} =
\num{390625}$ plantes, qui demandent \qty{39000}{m^2}, près de quatre
hectares. L'espace, la lumière et les nutriments, la faible portée d'un
stolon et la mort des plantes sous l'effet de l'encombrement arrêtent la
phase exponentielle en un ou deux ans~; le \omterm{def:b2:asexual-reproduction:clone}{clone} progresse alors comme
un front.
\end{solution}

\begin{solution}{exo:b2:asexual-reproduction:6}
$\num{14000}/130 \approx 110$ générations de troncs. Rayon d'un disque
de \qty{4.3e5}{m^2}~: $\sqrt{4.3\times 10^{5}/\pi} = \qty{370}{m}$~;
progression $370/\num{14000} = \qty{0.026}{m}$ par an, moins de trois
centimètres.
\end{solution}

\begin{solution}{exo:b2:asexual-reproduction:7}
$p_t = 2^{t}p_0/(1 + (2^{t} - 1)p_0)$~: $t = 5$~: $0.0032$~; $t = 10$~:
$0.1024/1.1023 = 0.093$~; $t = 15$~: $3.28/4.28 = 0.77$. Il ne se répand
jamais si sa valeur sélective est inférieure à la moitié de celle de la
lignée sexuée~: une valeur sélective relative $w$ donne $q' = 2wq$, qui
décroît pour $w <
\tfrac12$.
\end{solution}

\begin{solution}{exo:b2:asexual-reproduction:8}
$10 : 1$~: des racines~; $1 : 1$~: un cal indifférencié~; $1 : 10$~: des
tiges. Une bouture possède déjà une tige et a besoin de racines, que
l'auxine induit à la base coupée.
\end{solution}

\begin{solution}{exo:b2:asexual-reproduction:9}
1, 10, 100, \num{1000}, \num{10000} kg~: après la troisième saison,
seulement \qty{1}{t}, après la quatrième, \qty{10}{t}~: quatre saisons.
La vraie \omterm{def:b2:angiosperm-reproduction:seed}{graine} donne des milliers de \omterm{def:b2:angiosperm-reproduction:seed}{graines} par plante, mais la pomme
de terre est un tétraploïde \omterm{def:b2:meiosis-heredity:vocabulary}{hétérozygote}~: les plantules issues de
semis ségrègent, et aucune n'est la variété.
\end{solution}

\begin{solution}{exo:b2:asexual-reproduction:10}
$e^{-5} = 0.0067$~: 670 individus sans \omterm{def:b2:genome-diversification:mutation}{mutation} dans la grande
population, environ 7 dans la petite. Sept individus peuvent tous, par
hasard, ne laisser aucun descendant en quelques générations (la
probabilité qu'aucun ne se reproduise est de l'ordre de $e^{-7}$ par
génération), après quoi la meilleure classe est perdue à jamais et la
suivante devient la meilleure~: un cran du cliquet, répété jusqu'à ce
que le fardeau devienne insupportable. Avec 670, la perte est
pratiquement impossible --- jusqu'à ce qu'un goulot d'étranglement
réduise $N$ ou que le \omterm{thm:b2:genome-diversification:rate}{taux de mutation} augmente~; le cliquet ne tourne
jamais que dans un sens.
\end{solution}

\begin{solution}{exo:b2:asexual-reproduction:11}
Dans le \omterm{def:b2:asexual-reproduction:clone}{clone}, tous les hôtes sont sensibles~: la souche se propage sans
frein, à la vitesse à laquelle voyagent les spores. Dans la population
sexuée, le \omterm{def:b2:meiosis-heredity:vocabulary}{génotype} qu'elle contourne est l'une de $2^{10} = 1024$
combinaisons, présente chez environ un millième des plantes, dispersées
parmi des voisines résistantes~: l'épidémie s'essouffle faute d'hôtes.
La Lumper était un \omterm{def:b2:asexual-reproduction:clone}{clone} unique cultivé sur la majeure partie d'un
pays~; le mildiou trouva toutes les plantes sans défense, et il n'y eut
aucune variation à partir de laquelle sélectionner jusqu'à l'importation
d'autres variétés.
\end{solution}

\begin{solution}{exo:b2:asexual-reproduction:12}
Le \omterm{prop:b2:asexual-reproduction:whysex}{cliquet de Muller}, la combinaison des \omterm{def:b2:genome-diversification:mutation}{mutations} favorables et la
Reine rouge procurent chacun à la sexualité un bénéfice qui peut
dépasser un facteur deux dans un monde vaste, parasité et changeant. Les
rotifères bdelloïdes sont le cas embarrassant~: quarante millions
d'années sans mâles et sans effondrement par le cliquet. Ils semblent y
échapper par une résistance extrême à la dessiccation (qui tue leurs
parasites), par l'intégration d'ADN étranger et par une \omterm{prop:b2:genome-diversification:repair}{réparation de
l'ADN} d'une efficacité inhabituelle --- des exceptions qui montrent ce
qu'exige normalement la règle.
\end{solution}

\begin{solution}{pb:b2:asexual-reproduction:1}
\textbf{1.} $4\times 10^{4}$ plantes.
\textbf{2.} $1 + 6\times 2 = 13$.
\textbf{3.} \num{1300}~; \num{16900}~; \num{219700}.
\textbf{4.} $13^{t} = 400$~: $t = \ln 400/\ln 13 = 5.99/2.56 = 2.3$
ans --- le champ est plein au cours de la troisième année.
\textbf{5.} Rayon d'un disque d'un hectare~: $\sqrt{10^{4}/\pi} = \qty{56}{m}$~;
à \qty{0.5}{m} par an, \num{113} ans.
\textbf{6.} Cent fondateurs dispersés remplissent le champ de façon
exponentielle à partir de cent fronts en deux ans~; un seul fondateur
remplirait son voisinage en deux ans, puis progresserait lentement
pendant un siècle.
\textbf{7.} $5^{n} \ge 10^{6}$~: $n \ge 8.6$, soit neuf repiquages
($1.95\times 10^{6}$ tiges), 36~semaines.
\textbf{8.} $0.8\times 1.95\times 10^{6} = 1.6\times 10^{6}$ plantes.
\textbf{9.} $0.98$~; $0.98^{10} = 0.82$.
\textbf{10.} $10\times 0.98 = 9.8$ lignées saines attendues~; à partir
de dix boutures, une seule.
\textbf{11.} Le virus circule par les plasmodesmes et le phloème~; or le
dôme apical n'a pas de connexion vasculaire et se divise plus vite que
le virus ne se propage, si bien que le dixième de millimètre le plus
externe est en général indemne.
\textbf{12.} Laboratoire~: des plantes uniformes et indemnes de virus
d'emblée, mais coûteuses, et toute \omterm{def:b2:genome-diversification:mutation}{mutation} ou contamination de la
culture est multipliée un million de fois. Stolons~: gratuits et sur
place, mais trois saisons perdues et tous les virus des plantes mères
transmis.
\textbf{13.} Les apomictiques sont au nombre de $pN$ et donnent $F$
\omterm{def:b2:angiosperm-reproduction:seed}{graines} chacun~; les sexués, au nombre de $(1 - p)N$, donnent $F/2$
\omterm{def:b2:angiosperm-reproduction:seed}{graines} chacun (la moitié de leur effort va au \omterm{def:b2:plant-life-cycles:heterospory}{pollen}, qui ne produit
pas de \omterm{def:b2:angiosperm-reproduction:seed}{graines})~; la part des \omterm{def:b2:angiosperm-reproduction:seed}{graines} apomictiques vaut
$Fp/(Fp + F(1 - p)/2)
= 2p/(1 + p)$.
\textbf{14.} $p_5 = 0.032/1.031 = 0.031$~; $p_{10} = 1.024/2.023 =
0.51$~: il dépasse la moitié à la génération 10.
\textbf{15.} $p' = 2p(1 - cp)/\bigl(2p(1 - cp) + (1 - p)\bigr)$.
\textbf{16.} $p' = p$ lorsque $2(1 - cp^*) = 1$~: $p^* = 1/2c = 0.625$.
En dessous de $p^*$, $2(1 - cp) > 1$ et $p$ augmente~; au-dessus, il
diminue~: l'équilibre est stable.
\textbf{17.} $p^* \ge 1$ lorsque $c \le \tfrac12$~: l'apomictique envahit
la population, à moins que le parasite ne lui coûte plus de la moitié de
sa production lorsqu'il est commun.
\textbf{18.} $c = 0.6$~: $p^* = 0.83$. Pour maintenir la sexualité dans
la population, le parasite doit, à forte fréquence d'apomictiques,
réduire la valeur sélective du \omterm{def:b2:asexual-reproduction:clone}{clone} de plus que le double avantage ---
la Reine rouge doit courir vite.
\textbf{19.} Trois jeux de chromosomes ne peuvent pas s'apparier à la
\omterm{def:b2:meiosis-heredity:meiosis}{méiose}~: ses gamètes sont déséquilibrés~; il est exclu de la
recombinaison, et le \omterm{prop:b2:asexual-reproduction:whysex}{cliquet de Muller} et les parasites agissent sur lui
sans opposition --- ce qui explique que les lignées apomictiques de
pissenlits soient jeunes et sans cesse reconstituées à partir d'ancêtres
sexués.
\textbf{20.} $10^{-9}\times 4.8\times 10^{8} = 0.48$ par division.
\textbf{21.} $2\times \num{14000}\times 20 = 5.6\times 10^{5}$
divisions~; $0.48\times 5.6\times 10^{5} = 2.7\times 10^{5}$
\omterm{def:b2:genome-diversification:mutation}{mutations}.
\textbf{22.} $2.7\times 10^{5}/4.8\times 10^{8} = 5.6\times 10^{-4}$ du
génome~; environ \num{5400} différences codantes.
\textbf{23.} La compétition entre lignées cellulaires dans le méristème
(une lignée mutante qui se divise plus lentement est supplantée) et
entre \omterm{def:b2:asexual-reproduction:clone}{ramets} (un tronc porteur d'une \omterm{def:b2:genome-diversification:mutation}{mutation} défavorable drageonne
moins). Ce crible est plus faible car la plupart des \omterm{def:b2:genome-diversification:mutation}{mutations}
somatiques sont \omterm{def:b2:meiosis-heredity:vocabulary}{hétérozygotes} et masquées, et rien ne les expose~:
aucune \omterm{def:b2:meiosis-heredity:meiosis}{méiose} ne les rend \omterm{def:b2:meiosis-heredity:vocabulary}{homozygotes} et aucun stade zygote \omterm{def:b2:unicellular-diversity:unicellular}{unicellulaire}
ne teste le génome entier dans une seule cellule.
\textbf{24.} Il n'est pas uniforme~: ses troncs diffèrent en des
centaines de milliers de sites. Mais ces différences sont surtout
neutres et masquées, si bien que le \omterm{def:b2:asexual-reproduction:clone}{clone} est en pratique un seul
\omterm{def:b2:meiosis-heredity:vocabulary}{génotype} --- une mosaïque de variants somatiques d'un même \omterm{def:b2:asexual-reproduction:clone}{génet}.
\textbf{25.} Champ plein au bout de 2.3 ans~; un méristème donne
$1.6\times 10^{6}$ plantes en neuf mois~; avec $c = 0.8$, l'apomictique
se stabilise à $p^* = 0.625$, les sexués conservant \qty{37.5}{\%}.
\end{solution}
